\documentclass[11pt,a4paper]{article}

% =========================================================
% RVR-DRO-001
% İTÜ ROVER TAKIMI DRONE ALT EKİBİ
% DOKÜMAN HAZIRLAMA VE YAYIN STANDARDI
% =========================================================

% ---------------------------------------------------------
% PAKETLER
% ---------------------------------------------------------

\usepackage[utf8]{inputenc}
\usepackage[T1]{fontenc}
\usepackage[turkish]{babel}
\usepackage{lmodern}

\usepackage[
    top=22mm,
    bottom=22mm,
    left=20mm,
    right=20mm
]{geometry}

\usepackage{graphicx}
\usepackage{tabularx}
\usepackage{array}
\usepackage{booktabs}
\usepackage{longtable}
\usepackage{enumitem}
\usepackage{xcolor}
\usepackage{titlesec}
\usepackage{fancyhdr}
\usepackage{lastpage}
\usepackage{hyperref}
\usepackage{setspace}
\usepackage{tikz}
\usepackage{tcolorbox}
\usepackage{caption}
\usepackage{listings}
\usepackage{float}

% ---------------------------------------------------------
% RENKLER
% ---------------------------------------------------------

\definecolor{RoverKoyu}{HTML}{1B1B1B}
\definecolor{RoverGri}{HTML}{555555}
\definecolor{RoverAcikGri}{HTML}{E8E8E8}
\definecolor{RoverVurgu}{HTML}{B41F24}

% ---------------------------------------------------------
% PDF AYARLARI
% ---------------------------------------------------------

\hypersetup{
    colorlinks=true,
    linkcolor=RoverKoyu,
    urlcolor=RoverVurgu,
    pdftitle={RVR-DRO-001 Doküman Hazırlama ve Yayın Standardı},
    pdfauthor={İTÜ Rover Takımı Drone Alt Ekibi}
}

% ---------------------------------------------------------
% GENEL BİÇİM
% ---------------------------------------------------------

\setlength{\parindent}{0pt}
\setlength{\parskip}{6pt}
\renewcommand{\arraystretch}{1.3}

% ---------------------------------------------------------
% BAŞLIK BİÇİMLERİ
% ---------------------------------------------------------

\titleformat{\section}
  {\Large\bfseries\color{RoverKoyu}}
  {\thesection.}
  {0.7em}
  {}

\titleformat{\subsection}
  {\large\bfseries\color{RoverKoyu}}
  {\thesubsection.}
  {0.7em}
  {}

\titleformat{\subsubsection}
  {\normalsize\bfseries\color{RoverGri}}
  {\thesubsubsection.}
  {0.7em}
  {}

% ---------------------------------------------------------
% ŞEKİL VE TABLO
% ---------------------------------------------------------

\captionsetup{
    font=small,
    labelfont=bf
}

% ---------------------------------------------------------
% KOD GÖSTERİMİ
% ---------------------------------------------------------

\lstset{
    basicstyle=\ttfamily\small,
    frame=single,
    breaklines=true,
    columns=fullflexible,
    showstringspaces=false
}

% =========================================================
% BELGE BİLGİLERİ
% =========================================================

\newcommand{\BelgeNo}{RVR-DRO-001}
\newcommand{\BelgeAdi}{DOKÜMAN HAZIRLAMA VE YAYIN STANDARDI}
\newcommand{\Revizyon}{REV A}
\newcommand{\YayinTarihi}{22 EYLÜL 2026}
\newcommand{\BirimAdi}{İTÜ ROVER TAKIMI -- DRONE ALT EKİBİ}

% =========================================================
% ÜST BİLGİ / ALT BİLGİ
% =========================================================

\pagestyle{fancy}
\fancyhf{}

\fancyhead[L]{
    \begin{minipage}[c]{0.16\textwidth}
        \includegraphics[height=11mm]{itu_rover_logo.png}
    \end{minipage}
}

\fancyhead[C]{
    \begin{minipage}[c]{0.55\textwidth}
        \centering
        \small\bfseries \BirimAdi
    \end{minipage}
}

\fancyhead[R]{
    \begin{minipage}[c]{0.20\textwidth}
        \raggedleft
        \small\bfseries
        \BelgeNo\\
        \Revizyon
    \end{minipage}
}

\fancyfoot[L]{\scriptsize İTÜ ROVER TAKIMI}
\fancyfoot[C]{\scriptsize \Revizyon\ -- \YayinTarihi}
\fancyfoot[R]{\scriptsize SAYFA \thepage\ / \pageref{LastPage}}

\renewcommand{\headrulewidth}{0.8pt}
\renewcommand{\footrulewidth}{0.6pt}

\setlength{\headheight}{32pt}

% =========================================================
% ÖZEL KUTULAR
% =========================================================

\newtcolorbox{notkutusu}{
    colback=white,
    colframe=RoverGri,
    boxrule=0.7pt,
    arc=0mm,
    title=\textbf{NOT}
}

\newtcolorbox{dikkatkutusu}{
    colback=white,
    colframe=RoverVurgu,
    boxrule=1pt,
    arc=0mm,
    title=\textbf{DİKKAT}
}

\newtcolorbox{ornek}{
    colback=RoverAcikGri,
    colframe=RoverGri,
    boxrule=0.5pt,
    arc=0mm,
    title=\textbf{ÖRNEK}
}

% =========================================================
% BELGE
% =========================================================

\begin{document}

% =========================================================
% KAPAK
% =========================================================

\begin{titlepage}

    \thispagestyle{empty}

    \begin{center}

        \vspace*{10mm}

        \includegraphics[width=0.30\textwidth]{itu_rover_logo.png}

        \vspace{14mm}

        {\Large\bfseries İTÜ ROVER TAKIMI}

        \vspace{3mm}

        {\Large\bfseries DRONE ALT EKİBİ}

        \vspace{20mm}

        \rule{\textwidth}{1.2pt}

        \vspace{8mm}

        {\Huge\bfseries \BelgeAdi}

        \vspace{8mm}

        \rule{\textwidth}{1.2pt}

        \vfill

        \begin{tabular}{rl}

            \textbf{BELGE NO:} &
            \BelgeNo
            \\[2mm]

            \textbf{REVİZYON:} &
            \Revizyon
            \\[2mm]

            \textbf{YAYIN TARİHİ:} &
            \YayinTarihi
            \\[2mm]

            \textbf{BİRİM:} &
            \BirimAdi

        \end{tabular}

        \vfill

        {\small
        Bu belge, İTÜ Rover Takımı Drone Alt Ekibi bünyesinde
        oluşturulan teknik ve idari dokümanların hazırlanması,
        numaralandırılması, yayımlanması ve revize edilmesine
        ilişkin standartları tanımlar.
        }

        \vspace{14mm}

    \end{center}

\end{titlepage}

% =========================================================
% BELGE BİLGİLERİ
% =========================================================

\section*{BELGE BİLGİLERİ}

\begin{tabularx}{\textwidth}{|>{\bfseries}p{45mm}|X|}
\hline
Belge No & \BelgeNo \\
\hline
Belge Adı & \BelgeAdi \\
\hline
Birim & \BirimAdi \\
\hline
Revizyon & \Revizyon \\
\hline
Yayın Tarihi & \YayinTarihi \\
\hline
Hazırlayan & Drone Alt Ekibi \\
\hline
Kontrol Eden & \\
\hline
Onaylayan & \\
\hline
\end{tabularx}

\vspace{10mm}

% =========================================================
% REVİZYON GEÇMİŞİ
% =========================================================

\section*{REVİZYON GEÇMİŞİ}

\begin{tabularx}{\textwidth}{|c|c|X|X|}
\hline

\textbf{REVİZYON} &
\textbf{TARİH} &
\textbf{AÇIKLAMA} &
\textbf{HAZIRLAYAN}
\\
\hline

A &
22 EYLÜL 2026 &
İlk yayın &
Drone Alt Ekibi
\\
\hline

&
&
&
\\
\hline

&
&
&
\\
\hline

\end{tabularx}

\newpage

% =========================================================
% İÇİNDEKİLER
% =========================================================

\tableofcontents

\newpage

% =========================================================
% 1. GİRİŞ
% =========================================================

\section{GİRİŞ}

\subsection{Amaç}

Bu standardın amacı, İTÜ Rover Takımı Drone Alt Ekibi tarafından
hazırlanan tüm teknik ve idari dokümanların ortak bir yapıya sahip
olmasını sağlamaktır.

Dokümanların;

\begin{itemize}
    \item aynı takım kimliğini taşıması,
    \item kolay tanımlanabilmesi,
    \item sürümlerinin takip edilebilmesi,
    \item farklı ekip üyeleri tarafından aynı yöntemle hazırlanabilmesi,
    \item yıllar sonra dahi anlaşılabilir ve kullanılabilir olması,
    \item teknik bilginin ekip içerisinde kalıcı hâle getirilmesi
\end{itemize}

amaçlanır.

\subsection{Kapsam}

Bu standart, İTÜ Rover Takımı Drone Alt Ekibi adına hazırlanan;

\begin{itemize}
    \item teknik manuelleri,
    \item kullanım kılavuzlarını,
    \item bakım dokümanlarını,
    \item test prosedürlerini,
    \item test raporlarını,
    \item yazılım dokümantasyonlarını,
    \item elektronik ve aviyonik dokümanlarını,
    \item mekanik tasarım dokümanlarını,
    \item uçuş prosedürlerini,
    \item kontrol listelerini,
    \item eğitim dokümanlarını,
    \item mülakat ve değerlendirme dokümanlarını,
    \item sistem tanım dokümanlarını
\end{itemize}

kapsar.

\subsection{Uygulama}

Yeni oluşturulacak tüm RVR-DRO serisi belgelerin bu standarda uygun
hazırlanması esastır.

Mevcut belgeler, gerekli görüldüğünde yeni revizyon yayımlanarak
bu standarda uyarlanabilir.

% =========================================================
% 2. DOKÜMAN NUMARALANDIRMA SİSTEMİ
% =========================================================

\newpage

\section{DOKÜMAN NUMARALANDIRMA SİSTEMİ}

\subsection{Genel Yapı}

Drone Alt Ekibi dokümanları aşağıdaki biçimde numaralandırılır:

\begin{center}

\Large
\textbf{RVR-DRO-XXX}

\end{center}

Burada;

\begin{center}

\begin{tabularx}{0.85\textwidth}{|c|X|}
\hline
\textbf{Alan} & \textbf{Anlamı} \\
\hline
RVR & İTÜ Rover Takımı \\
\hline
DRO & Drone Alt Ekibi \\
\hline
XXX & Üç basamaklı benzersiz doküman numarası \\
\hline
\end{tabularx}

\end{center}

\subsection{Numara Tahsisi}

Her dokümana yalnızca bir kez numara tahsis edilir.

Tahsis edilen numara daha sonra farklı bir belge için kullanılamaz.

Bir dokümanın adı veya kapsamı değişse dahi aynı doküman devam ediyorsa
belge numarası korunur ve revizyon artırılır.

\begin{dikkatkutusu}
Silinen, yürürlükten kaldırılan veya artık kullanılmayan bir dokümanın
numarası başka bir dokümana verilmez.
\end{dikkatkutusu}

\subsection{Temel Dokümanlar}

Başlangıç dokümanları aşağıdaki şekilde tanımlanmıştır:

\begin{center}

\begin{tabularx}{0.95\textwidth}{|c|X|}
\hline
\textbf{Belge No} & \textbf{Belge Adı} \\
\hline
RVR-DRO-000 & Mülakat Soruları \\
\hline
RVR-DRO-001 & Doküman Hazırlama ve Yayın Standardı \\
\hline
\end{tabularx}

\end{center}

Yeni belgeler sırasıyla RVR-DRO-002,
RVR-DRO-003 ve devam eden numaralarla yayımlanır.

% =========================================================
% 3. REVİZYON SİSTEMİ
% =========================================================

\newpage

\section{REVİZYON SİSTEMİ}

\subsection{Revizyon Gösterimi}

Doküman revizyonları aşağıdaki biçimde gösterilir:

\begin{center}

\Large
\textbf{REV A}
\end{center}

İlk yayımlanan sürüm \textbf{REV A} olarak tanımlanır.

Daha sonraki sürümler alfabetik olarak ilerler:

\begin{center}

REV A $\rightarrow$ REV B $\rightarrow$ REV C $\rightarrow$ REV D

\end{center}

\subsection{Revizyon Gerektiren Değişiklikler}

Aşağıdaki değişiklikler yeni bir revizyon gerektirir:

\begin{itemize}
    \item teknik bilgilerin değiştirilmesi,
    \item prosedür değişiklikleri,
    \item yeni bölüm eklenmesi,
    \item bölüm kaldırılması,
    \item sistem mimarisinin değişmesi,
    \item kullanılan donanım veya yazılımın değişmesi,
    \item kritik hata düzeltmeleri,
    \item görev veya test yönteminin değiştirilmesi.
\end{itemize}

Yalnızca küçük yazım hatalarının düzeltilmesi veyahut tasarım değişikliği durumunda yeni revizyon
gerekliliğine ekip lideri karar verir.

\subsection{Revizyon Geçmişi}

Her dokümanda revizyon geçmişi tablosu bulunmalıdır.

Tabloda en az aşağıdaki bilgiler yer almalıdır:

\begin{itemize}
    \item revizyon kodu,
    \item yayın tarihi,
    \item yapılan değişikliğin kısa açıklaması,
    \item hazırlayan kişi veya birim.
\end{itemize}

% =========================================================
% 4. DOSYA ADLANDIRMA STANDARDI
% =========================================================

\newpage

\section{DOSYA ADLANDIRMA STANDARDI}

\subsection{Genel Biçim}

Elektronik dosyaların adı aşağıdaki biçimde oluşturulur:

\begin{center}

\texttt{BELGENO\_REV-X\_BELGE-ADI.pdf}

\end{center}

Örnek:

\begin{ornek}
\texttt{RVR-DRO-001\_REV-A\_DOKUMAN-HAZIRLAMA-VE-YAYIN-STANDARDI.pdf}
\end{ornek}

Kaynak LaTeX dosyası:

\begin{ornek}
\texttt{RVR-DRO-001\_REV-A\_DOKUMAN-HAZIRLAMA-VE-YAYIN-STANDARDI.tex}
\end{ornek}

\subsection{Dosya Adlarında Kullanılacak Karakterler}

Uyumluluk amacıyla dosya adlarında;

\begin{itemize}
    \item Türkçe karakter kullanılmaması,
    \item boşluk yerine tire veya alt çizgi kullanılması,
    \item belge numarasının dosya adının başında bulunması,
    \item revizyon kodunun dosya adına eklenmesi
\end{itemize}

önerilir.

% =========================================================
% 5. BELGE YAPISI
% =========================================================

\newpage

\section{BELGE YAPISI}

\subsection{Zorunlu Bölümler}

Standart bir RVR-DRO dokümanı mümkün olduğu ölçüde aşağıdaki sırayı
izlemelidir:

\begin{enumerate}
    \item Kapak,
    \item Belge Bilgileri,
    \item Revizyon Geçmişi,
    \item İçindekiler,
    \item Giriş,
    \item Ana Teknik İçerik,
    \item Ekler,
    \item Belge Sonu.
\end{enumerate}

Belgenin türüne göre bazı bölümler kullanılmayabilir.

\subsection{Kapak}

Kapakta en az aşağıdaki bilgiler bulunmalıdır:

\begin{itemize}
    \item İTÜ Rover logosu,
    \item İTÜ Rover Takımı,
    \item Drone Alt Ekibi,
    \item belge adı,
    \item belge numarası,
    \item revizyon,
    \item yayın tarihi,
    \item birim adı.
\end{itemize}

\subsection{Belge Bilgileri}

Belge bilgileri sayfasında en az;

\begin{itemize}
    \item belge numarası,
    \item belge adı,
    \item birim,
    \item revizyon,
    \item yayın tarihi,
    \item hazırlayan,
    \item kontrol eden,
    \item onaylayan
\end{itemize}

alanları bulunmalıdır.

\subsection{İçindekiler}

Beş sayfadan uzun dokümanlarda içindekiler bölümü bulunması önerilir.

Teknik manuellerde içindekiler bölümü zorunludur.

% =========================================================
% 6. SAYFA DÜZENİ
% =========================================================

\newpage

\section{SAYFA DÜZENİ}

\subsection{Kağıt Boyutu}

Tüm standart dokümanlar A4 boyutunda hazırlanır.

\subsection{Kenar Boşlukları}

Standart kenar boşlukları:

\begin{center}

\begin{tabular}{|l|c|}
\hline
Üst & 22 mm \\
\hline
Alt & 22 mm \\
\hline
Sol & 20 mm \\
\hline
Sağ & 20 mm \\
\hline
\end{tabular}

\end{center}

\subsection{Üst Bilgi}

Kapak hariç sayfaların üst bilgisinde;

\begin{itemize}
    \item sol tarafta İTÜ Rover logosu,
    \item ortada İTÜ ROVER TAKIMI -- DRONE ALT EKİBİ,
    \item sağ tarafta belge numarası ve revizyon
\end{itemize}

bulunur.

\subsection{Alt Bilgi}

Sayfa alt bilgisinde;

\begin{itemize}
    \item sol tarafta İTÜ ROVER TAKIMI,
    \item ortada revizyon ve yayın tarihi,
    \item sağ tarafta mevcut sayfa ve toplam sayfa sayısı
\end{itemize}

yer alır.

Örnek:

\begin{center}

\texttt{İTÜ ROVER TAKIMI \hspace{10mm} REV A -- 22 EYLÜL 2026
\hspace{10mm} SAYFA 12 / 36}

\end{center}

% =========================================================
% 7. YAZIM VE DİL STANDARDI
% =========================================================

\newpage

\section{YAZIM VE DİL STANDARDI}

\subsection{Belge Dili}

RVR-DRO serisi dokümanların ana dili Türkçedir.

Teknik gereklilik olmadığı sürece İngilizce ve Türkçe ifadelerin
aynı cümle içerisinde karışık kullanılmasından kaçınılmalıdır.



\subsection{Teknik Terimler}

Uluslararası kabul görmüş ve Türkçede yerleşik karşılığı bulunmayan
teknik ifadeler özgün biçimleriyle kullanılabilir.

Örnekler:

\begin{itemize}
    \item PID,
    \item ESC,
    \item GPS,
    \item LiPo,
    \item IMU,
    \item UART,
    \item SPI,
    \item I2C,
    \item MAVLink,
    \item ArUco.
\end{itemize}

Bir kısaltma ilk kez kullanıldığında gerekli görülürse açıklaması
verilmelidir.

\subsection{Anlatım}

Teknik metin;

\begin{itemize}
    \item açık,
    \item kısa,
    \item ölçülebilir,
    \item yoruma mümkün olduğunca kapalı,
    \item teknik olarak doğrulanabilir
\end{itemize}

olmalıdır.

\begin{dikkatkutusu}
``Yaklaşık iyi'', ``biraz yüksek'', ``fazla hızlı'' gibi ölçülemeyen
ifadeler kritik teknik prosedürlerde kullanılmamalıdır.
\end{dikkatkutusu}

Bunun yerine mümkün olduğunda sayısal sınırlar verilmelidir.

% =========================================================
% 8. BAŞLIKLAR VE NUMARALANDIRMA
% =========================================================

\newpage

\section{BAŞLIKLAR VE NUMARALANDIRMA}

Ana bölümler sayısal olarak numaralandırılır:

\begin{center}

1. GİRİŞ

2. SİSTEM TANIMI

3. DONANIM

4. YAZILIM

\end{center}

Alt başlıklar aşağıdaki şekilde devam eder:

\begin{center}

3.1. Uçuş Kontrolcüsü

3.2. Güç Sistemi

3.2.1. Batarya

3.2.2. Güç Dağıtımı

\end{center}

Aşırı derin başlık yapısından kaçınılmalıdır.

Mümkün olduğunca üç seviyeden daha fazla başlık kullanılmamalıdır.

% =========================================================
% 9. TABLOLAR
% =========================================================

\newpage

\section{TABLOLAR}

Her tablo;

\begin{itemize}
    \item anlaşılır bir başlığa,
    \item gerekli birimlere,
    \item açık sütun adlarına
\end{itemize}

sahip olmalıdır.

Teknik değerler mümkün olduğunca tablo hâlinde sunulmalıdır.

Örnek:

\begin{table}[H]

\centering

\caption{Örnek sistem parametreleri}

\begin{tabular}{|l|c|c|}
\hline
\textbf{Parametre} & \textbf{Değer} & \textbf{Birim} \\
\hline
Batarya gerilimi & 14.8 & V \\
\hline
Azami akım & 40 & A \\
\hline
Pervane çapı & 10 & inç \\
\hline
\end{tabular}

\end{table}

% =========================================================
% 10. ŞEKİLLER VE GÖRSELLER
% =========================================================

\newpage

\section{ŞEKİLLER VE GÖRSELLER}

Şekiller;

\begin{itemize}
    \item yeterli çözünürlükte,
    \item okunabilir,
    \item belge içerisinde gerekli boyutta,
    \item mümkünse beyaz veya şeffaf arka planlı
\end{itemize}

olmalıdır.

Her teknik şeklin altında açıklayıcı başlık bulunmalıdır.

Örnek:

\begin{verbatim}
\begin{figure}[H]
    \centering
    \includegraphics[width=0.75\textwidth]{drone_system.png}
    \caption{Drone sistemi genel yerleşimi}
\end{figure}
\end{verbatim}

Şekil üzerinde oklar,
numaralar ve açıklamalar kullanılıyorsa
metinler baskı veya PDF görüntüleme sırasında okunabilir olmalıdır.

% =========================================================
% 11. TEKNİK DEĞERLER VE BİRİMLER
% =========================================================

\newpage

\section{TEKNİK DEĞERLER VE BİRİMLER}

Teknik dokümanlarda mümkün olduğunca SI birim sistemi kullanılmalıdır.

Örnek birimler:

\begin{center}

\begin{tabular}{|l|c|}
\hline
Uzunluk & mm, m \\
\hline
Kütle & g, kg \\
\hline
Gerilim & V \\
\hline
Akım & A \\
\hline
Güç & W \\
\hline
Sıcaklık & $^\circ$C \\
\hline
Frekans & Hz \\
\hline
Hız & m/s \\
\hline
\end{tabular}

\end{center}

Bir sayısal değer ilk defa verildiğinde birimi belirtilmelidir.

Belirsiz veya ölçülmemiş değerler kesin değer gibi yazılmamalıdır.

% =========================================================
% 12. YAZILIM VE KOD BLOKLARI
% =========================================================

\newpage

\section{YAZILIM VE KOD BLOKLARI}

Kaynak kod,
terminal komutu,
yapılandırma dosyası
ve benzeri içerikler normal metinden ayrı gösterilmelidir.

Örnek:

\begin{lstlisting}
python3 main.py --live-low-alt
\end{lstlisting}

Terminal komutları mümkün olduğunca doğrudan kopyalanabilir biçimde
verilmelidir.

Bir komut belirli bir klasörde çalıştırılacaksa
çalışma dizini ayrıca belirtilmelidir.

Yapılandırma değerlerinde mümkün olduğunca;

\begin{itemize}
    \item parametre adı,
    \item varsayılan değer,
    \item kullanılan değer,
    \item birim,
    \item kısa açıklama
\end{itemize}

verilmelidir.

% =========================================================
% 13. NOT, DİKKAT VE UYARI KULLANIMI
% =========================================================

\newpage

\section{NOT, DİKKAT VE UYARI KULLANIMI}

Teknik bilgilerin önem seviyesini belirtmek için standart bilgi kutuları
kullanılabilir.

\subsection{Not}

Ek açıklama,
ipuçları
veya doğrudan güvenlik riski oluşturmayan bilgiler için kullanılır.

\begin{notkutusu}
Bu işlem sırasında uçuş kontrolcüsünün bilgisayara USB üzerinden bağlı
olması yeterlidir.
\end{notkutusu}

\subsection{Dikkat}

Hatalı uygulanması durumunda ekipman hasarına,
veri kaybına
veya görevin başarısız olmasına yol açabilecek işlemlerde kullanılır.

\begin{dikkatkutusu}
Pervane yönleri doğrulanmadan uçuş testi yapılmamalıdır.
\end{dikkatkutusu}

\subsection{Emniyet Uyarıları}

İnsan güvenliğini doğrudan ilgilendiren durumlarda uyarı açık,
kısa ve belirsizliğe yer vermeyecek şekilde yazılmalıdır.

% =========================================================
% 14. PROSEDÜRLERİN YAZILMASI
% =========================================================

\newpage

\section{PROSEDÜRLERİN YAZILMASI}

Bir işlemin sıralı olarak gerçekleştirilmesi gerekiyorsa
numaralandırılmış prosedür kullanılmalıdır.

Örnek:

\begin{enumerate}

    \item Batarya bağlantısını kontrol et.

    \item Pervanelerin doğru yönde takıldığını doğrula.

    \item Uçuş kontrolcüsünü enerjilendir.

    \item Yer kontrol istasyonu bağlantısını doğrula.

    \item Sensör durumlarını kontrol et.

    \item Uçuş alanının güvenli olduğunu doğrula.

\end{enumerate}

Her işlem maddesi mümkün olduğunca tek bir eylem içermelidir.

Bir adımın başarılı olup olmadığının anlaşılması için gerekli durumlarda
beklenen sonuç ayrıca yazılmalıdır.

% =========================================================
% 15. TEST DOKÜMANLARI
% =========================================================

\newpage

\section{TEST DOKÜMANLARI}

Test prosedürleri mümkün olduğunca aşağıdaki bilgileri içermelidir:

\begin{itemize}
    \item testin amacı,
    \item kullanılan sistem veya prototip,
    \item kullanılan donanımlar,
    \item yazılım sürümü,
    \item test ortamı,
    \item ön koşullar,
    \item test adımları,
    \item ölçülecek değerler,
    \item başarı kriterleri,
    \item test sonucu,
    \item gözlemler,
    \item ilgili log veya veri dosyaları.
\end{itemize}

Test gerçekleştirilmeden önce başarı kriterleri belirlenmelidir.

\begin{dikkatkutusu}
Test tamamlandıktan sonra sonuca göre başarı kriteri değiştirilmemelidir.
\end{dikkatkutusu}

% =========================================================
% 16. KAYNAK VE REFERANSLAR
% =========================================================

\newpage

\section{KAYNAK VE REFERANSLAR}

Başka bir RVR-DRO dokümanına atıf yapılırken belge numarası kullanılmalıdır.

Örnek:

\begin{ornek}
Uçuş öncesi kontroller için RVR-DRO-XXX dokümanına bakınız.
\end{ornek}

Harici teknik dokümanlarda mümkünse;

\begin{itemize}
    \item doküman adı,
    \item üretici veya yayımlayan kurum,
    \item sürüm,
    \item erişim tarihi
\end{itemize}

belirtilmelidir.

Üretici veri sayfaları,
resmî kullanım kılavuzları
ve teknik standartlar mümkün olduğunca öncelikli kaynak olarak
kullanılmalıdır.

% =========================================================
% 17. EKLER
% =========================================================

\newpage

\section{EKLER}

Ana belgenin akışını bozacak ölçüde uzun olan ancak belgeyle doğrudan
ilişkili içerikler ek bölümünde verilebilir.

Örnek ekler:

\begin{itemize}
    \item elektrik şemaları,
    \item kablolama tabloları,
    \item pin listeleri,
    \item parametre tabloları,
    \item parça listeleri,
    \item yazılım yapılandırmaları,
    \item uzun test sonuçları,
    \item kontrol listeleri.
\end{itemize}

Ekler aşağıdaki biçimde adlandırılır:

\begin{center}

EK A -- KABLOLAMA ŞEMASI

EK B -- PARAMETRE LİSTESİ

EK C -- KONTROL LİSTESİ

\end{center}

% =========================================================
% 18. DOKÜMAN HAZIRLAMA SÜRECİ
% =========================================================

\newpage

\section{DOKÜMAN HAZIRLAMA SÜRECİ}

Yeni bir doküman hazırlanırken aşağıdaki süreç uygulanır:

\begin{enumerate}

    \item Dokümana ihtiyaç olduğu belirlenir.

    \item Kullanılmamış bir RVR-DRO numarası tahsis edilir.

    \item Doküman adı belirlenir.

    \item RVR-DRO standart LaTeX şablonu kullanılır.

    \item Teknik içerik hazırlanır.

    \item İçerik teknik olarak kontrol edilir.

    \item Yazım ve biçim kontrolü yapılır.

    \item Doküman PDF olarak oluşturulur.

    \item Son PDF görsel olarak kontrol edilir.

    \item Revizyon tablosu tamamlanır.

    \item Doküman ilgili ekip depolama alanında yayımlanır.

\end{enumerate}

% =========================================================
% 19. REVİZYON VE DEĞİŞİKLİK SÜRECİ
% =========================================================

\newpage

\section{REVİZYON VE DEĞİŞİKLİK SÜRECİ}

Yayımlanmış bir belge değiştirilecekse mevcut dosya doğrudan üzerine
yazılmamalıdır.

Yeni revizyon oluşturulur.

Örneğin:

\begin{center}

RVR-DRO-005 REV A

$\downarrow$

RVR-DRO-005 REV B

\end{center}

Revizyon geçmişine yapılan değişiklik kısa ve anlaşılır şekilde yazılır.

Örnek:

\begin{center}

\begin{tabularx}{0.95\textwidth}{|c|c|X|}
\hline
REV & TARİH & AÇIKLAMA \\
\hline
A & 22.09.2026 & İlk yayın \\
\hline
B & 04.10.2026 & Uçuş öncesi kontrol prosedürü güncellendi \\
\hline
C & 18.10.2026 & Haberleşme sistemi bölümü eklendi \\
\hline
\end{tabularx}

\end{center}

% =========================================================
% 20. ARŞİVLEME
% =========================================================

\newpage

\section{ARŞİVLEME}

Yayımlanan dokümanların hem kaynak dosyası hem PDF çıktısı saklanmalıdır.

Önerilen klasör yapısı:

\begin{lstlisting}
DOCUMENTATION/
|
+-- RVR-DRO-000/
|   +-- SOURCE/
|   +-- RELEASE/
|   +-- ARCHIVE/
|
+-- RVR-DRO-001/
|   +-- SOURCE/
|   +-- RELEASE/
|   +-- ARCHIVE/
|
+-- RVR-DRO-002/
    +-- SOURCE/
    +-- RELEASE/
    +-- ARCHIVE/
\end{lstlisting}

\texttt{SOURCE} klasörü düzenlenebilir kaynak dosyalarını,

\texttt{RELEASE} klasörü güncel yayımlanmış PDF sürümünü,

\texttt{ARCHIVE} klasörü eski revizyonları

içerir.

Eski revizyonlar silinmemelidir.

% =========================================================
% 21. DOKÜMAN KONTROL LİSTESİ
% =========================================================

\newpage

\section{DOKÜMAN YAYIN KONTROL LİSTESİ}

Bir doküman yayımlanmadan önce aşağıdaki kontroller yapılmalıdır.

\begin{tabularx}{\textwidth}{|X|c|}
\hline
\textbf{Kontrol} & \textbf{Durum} \\
\hline
Belge numarası doğru mu? & $\square$ \\
\hline
Belge adı doğru mu? & $\square$ \\
\hline
Revizyon bilgisi doğru mu? & $\square$ \\
\hline
Yayın tarihi doğru mu? & $\square$ \\
\hline
İTÜ Rover logosu doğru mu? & $\square$ \\
\hline
Birim adı doğru mu? & $\square$ \\
\hline
Üst bilgi tüm sayfalarda doğru mu? & $\square$ \\
\hline
Alt bilgi tüm sayfalarda doğru mu? & $\square$ \\
\hline
Sayfa numaraları doğru mu? & $\square$ \\
\hline
İçindekiler güncel mi? & $\square$ \\
\hline
Revizyon geçmişi güncel mi? & $\square$ \\
\hline
Tablo ve şekiller okunabilir mi? & $\square$ \\
\hline
Teknik değerlerin birimleri mevcut mu? & $\square$ \\
\hline
Yazım hataları kontrol edildi mi? & $\square$ \\
\hline
Teknik bilgiler doğrulandı mı? & $\square$ \\
\hline
Dosya adı standarda uygun mu? & $\square$ \\
\hline
Kaynak dosya arşivlendi mi? & $\square$ \\
\hline
PDF son kez görsel olarak kontrol edildi mi? & $\square$ \\
\hline
\end{tabularx}

% =========================================================
% 22. STANDARTTAN SAPMALAR
% =========================================================

\newpage

\section{STANDARTTAN SAPMALAR}

Dokümanın kullanım amacı nedeniyle bu standardın bazı maddelerinin
uygulanmasının uygun olmadığı durumlarda biçimsel değişiklik yapılabilir.

Ancak aşağıdaki unsurlar mümkün olduğunca korunmalıdır:

\begin{itemize}
    \item RVR-DRO belge numarası,
    \item revizyon kodu,
    \item İTÜ Rover Takımı Drone Alt Ekibi tanımı,
    \item kurumsal üst bilgi,
    \item kurumsal alt bilgi,
    \item revizyon geçmişi,
    \item belge numarası ile dosya adının ilişkilendirilmesi.
\end{itemize}

Standarttan önemli ölçüde sapılan özel dokümanlarda gerekçe belge
içerisinde belirtilmelidir.

% =========================================================
% 23. SON HÜKÜMLER
% =========================================================

\newpage

\section{SON HÜKÜMLER}

Bu standart, İTÜ Rover Takımı Drone Alt Ekibi bünyesinde oluşturulan
RVR-DRO serisi dokümanlar için temel dokümantasyon referansıdır.

Yeni dokümanlar hazırlanırken bu belge esas alınmalıdır.

Dokümantasyon sisteminde yapılan yapısal değişiklikler öncelikle
bu dokümana işlenmeli ve yeni revizyon yayımlanmalıdır.

RVR-DRO serisinin biçimsel veya numaralandırma sistemini etkileyen
değişiklikler bireysel belgelerde bağımsız olarak uygulanmamalıdır.

% =========================================================
% BELGE SONU
% =========================================================

\newpage

\begin{center}

\vspace*{\fill}

{\Large\bfseries BELGE SONU}

\vspace{5mm}

{\bfseries \BelgeNo}

\vspace{2mm}

{\BelgeAdi}

\vspace{6mm}

{\small İTÜ Rover Takımı Drone Alt Ekibi}

\vspace*{\fill}

\end{center}

\end{document}